\begin{samepage}
\section{Research Questions}
\label{sec:intro-rq}

Five questions connect the implementation choices to the tests and remaining work:
\begin{enumerate}[label=\textbf{RQ\arabic*:}]
\item \textbf{Implementation.} Can the original charge-pump resonant circuit be built with the selected GaN half bridge, supplied planar magnetics, isolated sensing, auxiliary supplies, and USB-C output section?
\item \textbf{Control.} How is secondary voltage controlled during startup, loaded operation, load attachment, load removal, and no load, and how is that control carried from the first prototype to V2?
\item \textbf{Performance.} What output power, voltage regulation, input power quality, efficiency, switching behavior, and temperature are observed in the tests? How do those results compare with the \SI{65}{\watt} continuous minimum, the earlier \SI{75}{\watt} validation milestone, and the \SI{100}{\watt} continuous target?
\item \textbf{Magnetic integration.} How do the supplied inductor and transformer affect the connections, resonant design, mounting, insulation spacing, and thermal arrangement?
\item \textbf{Development limits.} What still needs to be measured or improved before the prototype can be rated as a USB-C power supply, and which V2/V3 objectives in the revised plan remain incomplete?
\end{enumerate}
\end{samepage}

The questions focus on the complete supply implemented and tested in this project. The PFC principle and the electromagnetic design of the supplied magnetic components came from earlier work. This project began with their integration into the USB-C supply.
